\subsection{What was already on disk}

Nothing in this study was generated for it. Every quantity comes from records
written by runs that finished before this analysis began, and each record is
bound to the manuscript by its SHA-256 digest; the figure code re-hashes each
file before reading it and stops the build on any mismatch. No image was
generated, no adapter was trained, and no evaluation was launched.

The records describe a subject-driven personalisation study on a latent
diffusion backbone. A concept is taught to the backbone through low-rank
adapters on the attention projections, and the resulting model is asked to
place that concept in a set of background scenes. Three arms exist on disk. The
\emph{jointly trained} arm puts one adapter on both attention types at once.
The \emph{composed} arm trains two adapters separately --- one on the
cross-attention projections, one on the self-attention projections --- and
switches both on at generation time. A third arm switches the self-attention
adapter off and leaves the cross-attention adapter alone.

\subsection{The comparison as it was recorded}

The registered comparison is between the jointly trained arm and the composed
arm on a single endpoint: the dispersion of identity fidelity across background
prompts. Each arm generated one image per prompt for \NPairs{} prompts, and the
dispersion is the standard deviation of the \NPairs{} resulting fidelities.
It is a population standard deviation, dividing by the number of prompts rather
than by one fewer; the figure code recomputes it from the per-prompt values and
refuses to proceed unless it reproduces the recorded number exactly, because
the entire comparison is between two such numbers.

The project recorded the comparison's evaluation pairing as
\texttt{\PairingLabel{}} and recorded that the pre-registered criterion
\KillFires{} on it. It also recorded, in the same file and in the ladder entry
that cites it, that this is an unpaired historical result rather than a paired
measurement. The present paper takes that caveat as its starting point and asks
what the same bytes do support.

\subsection{What the arms share}

Three configuration records, one per adapter, are bound here. The figure code
checks that the three agree on the backbone checkpoint, the training budget,
the learning rate, the adapter rank, the resolution, the training seed and the
instance prompt, and stops the build if any of them differ. It also checks
that the trainable parameter counts of the two composed adapters sum exactly to
the count the joint arm trains, which is the sense in which the two arms are
matched on capacity rather than merely similar.

\subsection{Reconstructing the pairs}

The per-prompt values underlying each dispersion are recorded, one row per
background prompt, with the prompt text as the row key. Rows are joined across
arms on that text rather than on position, so a reordering cannot silently pair
the wrong generations. Two endpoints are available per row: identity fidelity,
which the registered comparison aggregated, and prompt fidelity, which it did
not use.

A pair is one background prompt, and the difference is the composed arm's value
minus the jointly trained arm's value at that prompt. This recovers the
correspondence the dispersion discarded; it does not recover a pairing for the
registered endpoint itself, which cannot have one, because a dispersion across
prompts is a single number per arm.

\subsection{Tests, and what each of them can reach}

Three tests are reported on the same \NPairs{} pairs. The exact two-sided sign
test counts how many differences fall each way; ties would be dropped by
convention and are refused here instead, because with four pairs one discarded
tie changes the answer. The Wilcoxon signed-rank test on the same pairs is
reported beside it. The paired $t$ test is reported third, with its interval,
and is the only one of the three that assumes a distributional shape.

For the recorded comparison, which has no pairs, we report the interval its own
design supports: the ratio of the two dispersions with the confidence interval
an $F$ distribution gives at \NPairs{} generations per arm. That interval
assumes the two arms are independent, which they are not --- they share the
prompts. The assumption is stated rather than corrected because it is the
assumption the unpaired comparison implicitly makes, and the interval it yields
is already wide enough to settle the question this paper asks of it.

\subsection{The attainable probability of a test, before any data}

The exact sign test at $n$ pairs can return only the probabilities its
binomial reference distribution contains. The smallest two-sided value is
reached when every pair falls the same way and equals $2 \cdot 2^{-n}$; if $k$
pairs are allowed to fall the other way the smallest attainable value is
$2 \sum_{i \le k} \binom{n}{i} 2^{-n}$. These quantities depend on $n$ alone
and can be read before the experiment exists. They are what the pre-registered
design is sized against.

\subsection{How the protocol's size was chosen}

The protocol below is derived rather than picked. The background prompts are
the recorded grid and are not changed. The source design requires a second
concept, so that a conclusion about the mechanism can be distinguished from a
conclusion about one subject, and it requires the conclusion to be readable per
concept; each concept therefore needs enough pairs on its own. The number of
generation seeds is the smallest that leaves each concept able to reach
\Alpha{} with at least one pair falling the wrong way, which follows directly
from the attainable-probability expression above. Everything else in the
protocol table is a consequence of that choice.

\subsection{What is deliberately not computed}

The reference comparison this line of work would eventually be measured
against was never produced, and its tier is recorded as \BlockedTier{}. No
value is imputed for it, no substitute metric is introduced in its place, and
the two perceptual similarity measures the wider literature reports are
recorded on disk as not computed and are not computed here either. The recorded
gate list for the unrun lane stands at \AuthorisationState{}, and the paired
evaluation this paper specifies is not executed: the deliverable is the
specification.
